% ==============================================================================
\section{An Iterated-Greedy / VND Heuristic}
\label{sec:heuristic}
% ==============================================================================

The MILP of \S\ref{sec:milp} is exact but its size grows quickly: the
partition variables $z^-,\,z^+,\,b^B,\,b^C$ contribute
$O\!\left(|\calR|^2 \cdot |\calA| \cdot (N_{\max}+1)\right)$ binaries, so
proving optimality becomes intractable well before the fleet sizes of
operational interest (\S\ref{sec:experiments}).  We therefore develop a
metaheuristic, \textsc{IGVND}, that couples an Iterated-Greedy outer loop
with a Variable-Neighbourhood-Descent (VND) local search around a
problem-specific decoder.  The design draws on the NEH constructive
heuristic, Iterated Greedy, and Variable Neighbourhood
Search~\citep{Nawaz1983,Ruiz2007,Hansen2001,Mladenovic1997}.

% ------------------------------------------------------------------------------
\subsection{Overview}
\label{sec:igvnd_principle}
% ------------------------------------------------------------------------------

The central idea, recurrent in the scheduling literature the method draws
on, is to \emph{separate the expensive combinatorial decision from the
cheap timing decision}.  A solution is represented by a compact
\emph{combinatorial state} $(\pi, \sigma)$, where $\pi : \calR \to \calP$
is the position assignment and $\sigma$ is a global priority order over
the aircraft.  All search is conducted on this small state.  An inner,
deterministic \emph{decoder} maps a state $(\pi, \sigma)$ to concrete job
start and finish times, classifies every blocking access into its mode
(\S\ref{sec:problem}), and returns the objective value~\eqref{eq:milpobj};
we write $z(\pi,\sigma)$ for that value.  The decoder is the oracle that
prices a combinatorial state, and the search never manipulates timing
variables directly.

The method is built from five blocks, each detailed in a subsection below:
\begin{enumerate}[leftmargin=2.2em, itemsep=1pt]
  \item a \textbf{decoder} (\S\ref{sec:igvnd_decoder}) in two regimes ---
        a zero-movement regime that is feasible by construction, and a
        manoeuvre-aware regime that may spend Mode-(B)/Mode-(C) movements
        when they reduce the objective;
  \item a \textbf{construction scheme} (\S\ref{sec:igvnd_construction})
        built on two deterministic priority rules plus a rank-biased
        randomised shuffle that gives every further restart a distinct
        seed;
  \item a \textbf{Variable Neighbourhood Descent}
        (\S\ref{sec:igvnd_vnd}) that refines a state to a local optimum;
  \item an \textbf{Iterated-Greedy} loop (\S\ref{sec:igvnd_ig}) that
        perturbs and re-optimises around the incumbent; and
  \item a \textbf{multi-start} driver with a feasibility safety net
        (\S\ref{sec:igvnd_safety}).
\end{enumerate}

Algorithm~\ref{alg:igvnd} states the method at a high level.  Restarts are
drawn \emph{until the time budget is exhausted}: each restart takes a seed
order (deterministic for the first two, biased-randomised afterwards),
optimises it first in the zero-movement regime to obtain a
guaranteed-feasible \emph{floor}, and then polishes it in the
manoeuvre-aware regime, retaining the polished schedule only if an
independent feasibility check certifies it and it improves the floor.  A
single dense-nesting candidate is added when the weights reward it.  The
best certified schedule over all restarts is returned.  Because every
restart ends early once its own search stalls, small instances perform
hundreds of independent restarts within the budget while large ones
naturally receive few, without any instance-dependent tuning.

\begin{algorithm}[htbp]
\DontPrintSemicolon
\SetKwInput{KwData}{Input}
\SetKwInput{KwResult}{Output}
\KwData{instance $I$; weights $(W^M, W^D, W^S)$; time budget $T$.}
\KwResult{a feasible schedule.}
preprocess $I$; let $\rho_{\mathrm{NEH}}, \rho_{\mathrm{SLACK}}$ be the two
   deterministic priority rules (\S\ref{sec:igvnd_construction})\;
set the per-restart time cap $\tau = T/8$ if $R \le 10$, $\;T/4$ if
   $R \le 20$, $\;T/3$ otherwise\;
$z^\star \leftarrow +\infty$;\quad $i \leftarrow 1$\;
\While{time remains}{
  allot $\min(\tau,\ \text{remaining time})$ to restart $i$\;
  $\rho \leftarrow \rho_{\mathrm{NEH}}$ if $i{=}1$;\;\quad
     $\rho_{\mathrm{SLACK}}$ if $i{=}2$;\;\quad
     $\textsc{BiasedShuffle}$ of the better-scoring of the two if $i \ge 3$
     \tcp*{Alg.~\ref{alg:construct}}
  $(\pi, \sigma) \leftarrow \textsc{Construct}(\rho)$\;
  $(\pi, \sigma) \leftarrow \textsc{Search}(\pi, \sigma, \textsc{Decode}_0)$;
     \quad $z_{\mathrm{floor}} \leftarrow z_0(\pi, \sigma)$
     \tcp*{zero-movement floor}
  $(\pi', \sigma') \leftarrow \textsc{Search}(\pi, \sigma, \textsc{Decode}_1)$
     \tcp*{manoeuvre-aware polish}
  \eIf{$z_1(\pi', \sigma') < z_{\mathrm{floor}}$ \textbf{\emph{and}}
        $\textsc{Decode}_1(\pi', \sigma')$ passes the feasibility check}{
    $(\sigma^\circ, z^\circ) \leftarrow$ the polished schedule and its value\;
  }{
    $(\sigma^\circ, z^\circ) \leftarrow$ the floor schedule and $z_{\mathrm{floor}}$\;
  }
  \lIf{$z^\circ < z^\star$}{$z^\star \leftarrow z^\circ$; record the schedule}
  $i \leftarrow i + 1$\;
}
\If{$W^S$ dominates the weights \textbf{\emph{and}} $\calA \ne \emptyset$}{
  build the dense-nesting candidate (\S\ref{sec:igvnd_safety})\;
  adopt it if it is feasible and improves $z^\star$\;
}
\Return the recorded schedule attaining $z^\star$\;
\caption{\textsc{IGVND} --- master procedure.}
\label{alg:igvnd}
\end{algorithm}

% ------------------------------------------------------------------------------
\subsection{The decoder}
\label{sec:igvnd_decoder}
% ------------------------------------------------------------------------------

\paragraph{Zero-movement regime.}
This regime produces schedules that are feasible by construction with
\emph{no manoeuvres} ($n = 0$).  For every blocking arc it keeps each rear
aircraft's two access instants in Mode~A.  Relative to a front aircraft,
this leaves exactly three admissible placements of a rear aircraft's stay:
entirely \emph{before} the front, entirely \emph{after} it, or
\emph{enclosing} it (the rear enters at least $\eta$ before the front
starts and leaves at least $\eta$ after it finishes).  The third placement
--- \emph{nesting} --- is what makes the regime competitive: a long
aircraft can wrap a shorter one so that two blocking-related aircraft
overlap in time instead of being serialised.  Aircraft sharing a position
are separated by $\varepsilon$; aircraft on non-conflicting positions run
in parallel.  Processing the aircraft in the priority order $\sigma$, the
decoder places each one at the earliest start that respects every
constraint against the already-placed aircraft --- a forbidden-interval
scan that naturally lands in the hole left by the nesting option
(Algorithm~\ref{alg:decode0}).  Jobs are then packed tight, so
$\kappa_j = 0$ for every job and no Mode-(B)/Mode-(C) feedback on timing
can arise.  This regime is fast, always feasible, and already optimal on
instances with little or no blocking.

\begin{algorithm}[htbp]
\DontPrintSemicolon
\SetKwInput{KwData}{Input}
\SetKwInput{KwResult}{Output}
\KwData{state $(\pi, \sigma)$.}
\KwResult{a feasible schedule with $n = 0$, and its value $z_0$.}
$\mathcal{Q} \leftarrow \emptyset$ \tcp*{aircraft already placed}
\ForEach{$r \in \calR$ \emph{in the order} $\sigma$}{
  $\mathcal{F} \leftarrow \emptyset$ \tcp*{forbidden start instants for $r$}
  \ForEach{$r' \in \mathcal{Q}$}{
    \lIf{$\pi(r') = \pi(r)$}{add to $\mathcal{F}$ the band that violates the
         $\varepsilon$ separation from $r'$}
    \ForEach{$(p, p') \in \calA$ \emph{with} $\{\pi(r), \pi(r')\} = \{p, p'\}$}{
      add to $\mathcal{F}$ the bands that would place an access of the rear
         aircraft within $\eta$ of the front stay\;
      \tcp*[l]{the admissible residue is: before, after, or enclosing}
    }
  }
  $s_r \leftarrow \min\{\, t \ge E_r : t \notin \mathcal{F} \,\}$\;
  lay the jobs of $r$ back-to-back from $s_r$ with $\kappa_j = 0$\;
  $\mathcal{Q} \leftarrow \mathcal{Q} \cup \{r\}$\;
}
\Return the schedule and $z_0 = W^M m + W^D \sum_r v^D_r$ \tcp*{$n = 0$}
\caption{$\textsc{Decode}_0$ --- zero-movement decoder.}
\label{alg:decode0}
\end{algorithm}

\paragraph{Manoeuvre-aware regime.}
When the weights make manoeuvres worth their cost, the schedule can be
compressed further by \emph{spending} them.  The positions are processed
deepest-rear first, so that when a front aircraft is laid out the access
instants of all the rear aircraft it blocks are already fixed.  A front
aircraft's start is then chosen to minimise its \emph{local} contribution
to the objective over a candidate set that always contains the three
zero-movement placements as well as starts that align a job interior over
an access (inviting a cheap Mode~C) or a job end just before an access
(inviting a cheap Mode~B); the local cost of a start is evaluated by a
forward sweep that classifies every exposed access and prices the
resulting manoeuvres (Algorithm~\ref{alg:decode1}).  Because the
zero-movement placements remain candidates, this regime can only ever
\emph{add} the manoeuvre option --- it never discards a feasible
no-movement schedule.

\begin{algorithm}[htbp]
\DontPrintSemicolon
\SetKwInput{KwData}{Input}
\SetKwInput{KwResult}{Output}
\SetKwProg{Fn}{Procedure}{}{}
\SetKwFunction{FPlace}{PlaceFront}
\KwData{state $(\pi, \sigma)$.}
\KwResult{a feasible schedule, possibly with manoeuvres, and its value $z_1$.}
\ForEach{position $p \in \calP$ \emph{in order of decreasing blocking depth}}{
  $\phi \leftarrow -\infty$ \tcp*{finish of the previous aircraft at $p$}
  \ForEach{$r$ \emph{assigned to} $p$, \emph{in the order} $\sigma$}{
    $\ell \leftarrow \max\{E_r,\ \phi + \varepsilon\}$\;
    $A \leftarrow$ the access instants of every already-placed rear of $r$\;
    $(s_r, f_r) \leftarrow \FPlace{$r, \ell, A$}$;\quad $\phi \leftarrow f_r$\;
  }
}
\Return the schedule and its value $z_1$\;
\BlankLine
\Fn{\FPlace{$r$, $\ell$, $A$}}{
  $\mathcal{S} \leftarrow \{\ell\}
     \cup \{\, \tau - \eta - T_r,\ \tau - \eta,\ \tau + \eta : \tau \in A \,\}$
     \tcp*{before / after / nest}
  $\mathcal{S} \leftarrow \mathcal{S}
     \cup \{\text{starts aligning a job interior over some } \tau \in A\}$
     \tcp*{invite Mode C}
  $\mathcal{S} \leftarrow \mathcal{S}
     \cup \{\text{starts aligning a job end just before some } \tau \in A\}$
     \tcp*{invite Mode B}
  $c^\star \leftarrow +\infty$\;
  \ForEach{$s \in \mathcal{S}$ \emph{with} $s \ge \ell$}{
    sweep the jobs of $r$ from $s$, classifying each $\tau \in A$ as
       Mode~A, Mode~B (open a gap of width $\ge \mu$) or Mode~C
       (extend the job by $\delta$); reject $s$ if any $\tau$ cannot be
       classified or falls within $\eta$ of a job edge\;
    \If{$s$ \emph{is feasible}}{
      $c \leftarrow W^M f_r + W^D \max\{0, f_r - L_r\} + W^S \cdot 2\nu_s$
         \tcp*{$\nu_s$ = manoeuvres at $s$}
      \lIf{$c < c^\star$}{$c^\star \leftarrow c$; keep $s$ as the best start}
    }
  }
  \Return the best start and its finish\;
}
\caption{$\textsc{Decode}_1$ --- manoeuvre-aware decoder.}
\label{alg:decode1}
\end{algorithm}

% ------------------------------------------------------------------------------
\subsection{Construction: two rules and a biased shuffle}
\label{sec:igvnd_construction}
% ------------------------------------------------------------------------------

Each restart builds its seed by greedy insertion under a priority order
$\rho$, so restarts differ by construction and not only by the
perturbation source.  Two deterministic orders cover the objective's two
combinatorial levers: $-T_r$ (the NEH makespan rule --- the bulkiest
aircraft are placed while the schedule is still empty) and the slack
$L_r - E_r - T_r$ (least due-date headroom first, which steers
tight-target aircraft into early slots --- the delay lever).  Every
restart after the second draws its order from a single diversification
mechanism: a \emph{rank-biased geometric shuffle} of whichever
deterministic order scored better, in which the next aircraft is drawn
from the remaining base order with $\Pr(\text{rank } k) \propto
(1-\beta)^k$ (we use $\beta = 0.3$), so each restart explores a distinct
but still rule-shaped insertion order
(Algorithm~\ref{alg:construct}).  Combined with the open-ended restart
loop of Algorithm~\ref{alg:igvnd}, this gives the search an unlimited
supply of diverse seeds.  The cost of a tentative insertion is the value
returned by the zero-movement decoder restricted to the aircraft committed
so far.

\begin{algorithm}[htbp]
\DontPrintSemicolon
\SetKwInput{KwData}{Input}
\SetKwInput{KwResult}{Output}
\KwData{a priority order $\rho$; for the shuffle, a base order and the RNG.}
\KwResult{a seed state $(\pi, \sigma)$.}
\textbf{Greedy insertion under $\rho$:}\;
\Indp
  $\pi \leftarrow \emptyset$\;
  \ForEach{$r \in \calR$ \emph{in the order} $\rho$}{
    $\pi(r) \leftarrow \argmin_{p \in \calP}
       z_0\big(\pi \cup \{r \mapsto p\}\big)$
       \tcp*{cheapest position given the committed prefix}
  }
  \Return $(\pi,\ \rho)$\;
\Indm
\BlankLine
\textbf{\textsc{BiasedShuffle} of a base order $\rho_{\mathrm{base}}$:}\;
\Indp
  $\mathit{pool} \leftarrow \rho_{\mathrm{base}}$;\quad $\rho \leftarrow ()$\;
  \While{$\mathit{pool}$ is not empty}{
    draw a rank $k$ with $\Pr(k) \propto (1-\beta)^{k}$
       \tcp*{$\beta = 0.3$: most mass on ranks 0--3}
    move the $\min(k, |\mathit{pool}|-1)$-th element of $\mathit{pool}$
       to the end of $\rho$\;
  }
  \Return $\rho$ \tcp*{then construct greedily under $\rho$}
\Indm
\caption{\textsc{Construct} --- greedy insertion and the biased shuffle.}
\label{alg:construct}
\end{algorithm}

% ------------------------------------------------------------------------------
\subsection{Variable Neighbourhood Descent}
\label{sec:igvnd_vnd}
% ------------------------------------------------------------------------------

A seed is refined to a local optimum by a sequential VND with the basic
reset rule: the neighbourhoods are explored in a fixed sequence and, on any
improving move, the search returns to the first neighbourhood
(Algorithm~\ref{alg:vnd}).  Three neighbourhoods act on the combinatorial
state and delegate all timing to the decoder: $N_1$ \emph{reassign} (move
one aircraft to a different position), $N_2$ \emph{swap} (exchange the
positions of two aircraft), and $N_3$ \emph{reorder} (swap two aircraft in
the priority order $\sigma$).  Each neighbourhood is explored
first-improvement, so the first strictly improving move is taken
immediately.

\begin{algorithm}[htbp]
\DontPrintSemicolon
\SetKwInput{KwData}{Input}
\SetKwInput{KwResult}{Output}
\KwData{state $(\pi, \sigma)$; decoder $\textsc{Decode}_d$.}
\KwResult{a state that is locally optimal for $N_1, N_2, N_3$.}
$k \leftarrow 1$\;
\While{$k \le 3$}{
  scan the moves of $N_k$ and apply the first one that strictly lowers
     $z_d(\pi, \sigma)$\;
  \eIf{an improving move was applied}{
    $k \leftarrow 1$ \tcp*{reset to the first neighbourhood}
  }{
    $k \leftarrow k + 1$ \tcp*{advance to the next neighbourhood}
  }
}
\Return $(\pi, \sigma)$\;
\caption{\textsc{Vnd} --- sequential variable neighbourhood descent.}
\label{alg:vnd}
\end{algorithm}

% ------------------------------------------------------------------------------
\subsection{Iterated-Greedy outer loop}
\label{sec:igvnd_ig}
% ------------------------------------------------------------------------------

Around the VND runs an Iterated-Greedy loop (Algorithm~\ref{alg:search}).
Each iteration \emph{perturbs} the current state by removing the $k$
aircraft of largest contribution to the objective --- a delay-weighted
score with light randomisation, so the loop explores rather than repeating
the same removals --- and \emph{reconstructs} the state by reinserting each
removed aircraft at the position and order-slot that minimises the decoded
objective.  The perturbed state is re-optimised by the VND and accepted as
the new current state when it does not worsen the walk; the best state is
tracked separately, and the walk is restarted from it after a streak of
non-improving iterations.  The procedure is decoder-agnostic: the master
procedure calls it once with $\textsc{Decode}_0$ for the floor and once
with $\textsc{Decode}_1$ for the polish.

\begin{algorithm}[htbp]
\DontPrintSemicolon
\SetKwInput{KwData}{Input}
\SetKwInput{KwResult}{Output}
\SetKwProg{Fn}{Procedure}{}{}
\SetKwFunction{FPert}{Perturb}
\KwData{seed $(\pi, \sigma)$; decoder $\textsc{Decode}_d$; deadline.}
\KwResult{the best state found.}
$(\pi, \sigma) \leftarrow \textsc{Vnd}(\pi, \sigma, \textsc{Decode}_d)$\;
\emph{best} $\leftarrow$ \emph{cur} $\leftarrow (\pi, \sigma)$;\quad
   $t \leftarrow 0$ \tcp*{$t$ = non-improving streak}
\Repeat{the deadline is reached \emph{or} $t \ge t_{\max}$}{
  \emph{cand} $\leftarrow \textsc{Vnd}\big(\FPert{\emph{cur}},
     \textsc{Decode}_d\big)$\;
  \lIf{$z_d(\emph{cand}) \le z_d(\emph{cur})$}{\emph{cur} $\leftarrow$ \emph{cand}}
  \eIf{$z_d(\emph{cand}) < z_d(\emph{best})$}{
    \emph{best} $\leftarrow$ \emph{cand};\quad $t \leftarrow 0$\;
  }{
    $t \leftarrow t + 1$;\quad
    \textbf{if} $t$ is a multiple of the restart period \textbf{then}
       \emph{cur} $\leftarrow$ \emph{best}\;
  }
}
\Return \emph{best}\;
\BlankLine
\Fn{\FPert{$(\pi, \sigma)$}}{
  remove the $k$ aircraft of largest score
     $W^D v^D_r + \omega\, T_r$ (lightly randomised)\;
  \ForEach{removed aircraft $r$}{
    reinsert $r$ at the $(\text{position}, \text{slot in } \sigma)$
       minimising the decoded objective\;
  }
  \Return the rebuilt state\;
}
\caption{\textsc{Search} --- Iterated-Greedy loop around the VND.}
\label{alg:search}
\end{algorithm}

% ------------------------------------------------------------------------------
\subsection{Multi-start and the feasibility safety net}
\label{sec:igvnd_safety}
% ------------------------------------------------------------------------------

The construct-and-improve procedure is run from several seeds, and the best
result is kept.  The number of restarts is adaptive to instance size (8 for
$R \le 10$, 4 for $R \le 20$, 3 otherwise): more restarts on the cheap
small instances, fewer on the large ones that need their per-start time.
Because the search is time-limited and hence non-deterministic, a single
run occasionally settles in a poor high-delay basin, and independent
restarts make finding the good basin reliable.

Two safeguards make the method robust.  First, within each restart the
manoeuvre-aware polish is adopted only if an \emph{independent compliance
check} certifies it and it strictly improves the zero-movement floor
(Algorithm~\ref{alg:igvnd}); the method can therefore never return an
infeasible schedule and can never do worse than its zero-movement result,
whatever approximation the manoeuvre-aware decoder makes.  Second, when the
movement-priority profile meets a dense blocking topology, one further
candidate is built once and folded in by the same best-of / check rule: a
\emph{dense concentric-nesting} schedule written with explicit start times
rather than via the earliest-feasible decode.  In a dense component the
zero-movement optimum packs aircraft into a few concentric-nesting waves
--- a long aircraft wraps shorter ones, all overlapping within a wave's
span --- but the earliest-feasible decode cannot produce that, because it
always prefers placing \emph{before} over \emph{nested}.  The dedicated
builder sorts aircraft by stay length, groups them into waves of at most
$|\calP|$, makes the longest of each wave the outer container on the
deepest position, and \emph{stretches} shorter aircraft with idle time so
the stay lengths step down by at least $2\eta$ (the nesting condition) even
when their work durations tie; a small beam over two wave partitions is
tried and the best kept.  On other topologies this candidate is dominated
and ignored by the best-of.
